% \sectionheader{Cantos de Misa}

\songsection{B - Cantos de Misa}
\hrule \vspace{0.5cm}

\renewcommand{\songsectionletter}{B} % <-- Establece el prefijo T (texto plano)
\setcounter{songnum}{1}              % Reinicia el número a 1 (Resultado: T1

\begin{intersong}
\vspace{0.3cm}
\hspace*{-0.7cm}\textbf{\huge\textcolor{darkorange}{--- Entrada ---}}
\vspace{0.3cm}
\end{intersong}

% -----------------------------------------------------------------
\beginsong{Venimos hoy a tu altar}[by={Carlos Bravo} ]

\beginchorus
Ve\[G]nimos hoy a tu al\[Bm]tar,
a can\[C]tarte Se\[D]ñor,
pues \[C]tú eres la ale\[G]gría
\[D7]de nuestro cora\[G]zón.
pues \[C]tú eres la ale\[G]gría
\[D7]de nuestro cora\[G]zón. 
\endchorus

\beginverse
\[G]Tú hiciste los \[Em]cielos,
los llenas de es\[C]trellas, de luz y ca\[D7]lor.
\[G]Tú pintaste la a\[Em]urora,
hiciste las \[C]nubes, las puestas del \[D]sol.
\endverse

\beginverse
\chordsoff
Tú creaste la risa,
la paz y la dicha, la felicidad.
Tú al darnos la vida,
nos das tus riquezas, tu eterna amistad.
\endverse

\beginverse
\chordsoff
Tú nos diste a tu madre,
nos diste tu Cuerpo, tu Sangre en manjar.
Tú nos diste esperanza,
la fe y nos hiciste capaces de amar.
\endverse
\endsong

% ----------------------------------  STP ------------------------------------------------- %

\nextcol
\begin{intersong}
\vspace{0.3cm}
\hspace*{-0.7cm}\textbf{\huge\textcolor{darkorange}{--- Kyrie ---}}
\vspace{0.3cm}
\end{intersong}

% -----------------------------------------------------------------
\beginsong{Kyrie Eleison}[by={} ]

\beginchorus
\[Em]Kyrie \[C]Kyrie\[D]lei\[Em]son,
\[Em]Kyrie \[C]Kyrie\[D]lei\[Em]son
\endchorus

\beginverse
\[Em]Christe \[C]Christe-e\[D]lei\[Em]son,
\[Em]Christe \[C]Christe-e\[D]lei\[Em]son
\endverse

\beginchorus
\[Em]Kyrie \[C]Kyrie\[D]lei\[Em]son,
\[Em]Kyrie \[C]Kyrie\[D]lei\[Em]son
\endchorus
\endsong

% -------------------------------------------------------
\beginsong{Señor ten piedad (Jésed - Cuaresma)}[by={Federico Carranza (Jésed)}, sr={Disco: ``La cena del Señor Vol. I''}]

\beginchorus
{\nolyrics Intro: \[Em]}
Se\[C]ñor \[D]ten pie\[Em]dad
Se\[C]ñor \[D]ten pie\[Em]dad
\endchorus

\beginverse 
\[Am]Cristo \[G]ten pie\[D]dad
\[Am]Cristo \[G]ten pie\[D]dad
\endverse

\beginchorus
Se\[C]ñor \[D]ten pie\[Em]dad
Se\[C]ñor \[D]ten pie\[Em]dad
\endchorus

\iflink
  \textnote{\href{https://www.youtube.com/watch?v=nEkBYEsn7eA}{Consultar}} 
\fi
\endsong

% ----------------------------------  Gloria ------------------------------------------------- %
\nextcol
\begin{intersong}
\vspace{0.3cm}
\hspace*{-0.7cm}\textbf{\huge\textcolor{darkorange}{--- Gloria ---}}
\vspace{0.3cm}
\end{intersong}

% -------------------------------------------------------
\beginsong{Gloria a Dios}[by={}]
\vspace{-0.3cm}
\beginverse
{\nolyrics Intro: \[E] \[A] \[E]}
\endverse
\vspace{-0.3cm}

\beginchorus
\[E]Gloria a Dios en el cielo
y \[C#m]paz en la tierra
a los \[F#m]hombres que \[A]ama el Se\[B7]ñor
\endchorus

\beginverse
\[E]Por tu inmensa gloria, te ala\[C#m]bamos, bendecimos,
Te ado\[F#m]ramos y glo\[A]rificamos, Se\[B7]ñor.
\endverse

\beginverse 
Te damos \[A]gracias, Señor \[B7]Dios, Rey celes\[C#m]tial,
Dios \[A]Padre todopode\[B7]roso, Se\[C#m]ñor
\[A]Hijo \[E]único Jesu\[B7]cristo, Jesu\[A]cristo.
\endverse

\beginverse 
Se\[E]ñor Dios, Cordero, de \[C#m]Dios Hijo del Padre
Tú que \[F#m]quitas el pe\[A]cado ten pie\[B7]dad
A\[E]tiende a nuestra súplica
Tú que es\[C#m]tás a la derecha del \[F#m]Padre.
De no\[A]sotros ten pie\[B7]dad.
\endverse

\beginverse 
Porque sola\[A]mente Tú eres \[B7]Santo, sólo \[C#m]Tú,
Jesu\[A]cristo, el Al\[B7]tísimo Se\[C#m]ñor
con el \[A]Santo Espí\[E]ritu
En la \[B7]gloria de Dios \[A]Padre. A\[E]mén
\endverse

\endsong


% \begin{intersong}
% \vspace{0.3cm}
% \hspace*{-0.7cm}\textbf{\huge\textcolor{darkorange}{--- Aclamación ---}}
% \vspace{0.3cm}
% \end{intersong}


% ----------------------------------  Ofertorio ------------------------------------------------- %
\nextcol
\begin{intersong}
\vspace{0.3cm}
\hspace*{-0.7cm}\textbf{\huge\textcolor{darkorange}{--- Ofertorio ---}}
\vspace{0.3cm}
\end{intersong}


% -------------------------------------------------------
\beginsong{Con nuestras manos ya cansadas}[by={Benito Ortiz}]

\beginverse
\[Bm] Con nuestras \[A]manos ya can\[Bm]sadas, \[B7-G]
y nuestros \[A]pies de tanto an\[D]dar \[F#m-G]
Venimos \[A]hoy a presen\[D]tar\[F#m]nos  Se\[Bm]ñor, 
\[G]ante \[E7]tu al\[A]tar.
\endverse

\beginverse \chordsoff
Es nuestra ofrenda nuestro canto  
una oración universal 
Que te queremos ofrecer hoy  Señor 
Para que nos puedas perdonar. 
\endverse

\beginchorus
Y \[D]así, con el vino y \[A]con el pan
nuestras \[G]vidas \[Em] puedas cam\[A7]biar
Y \[D]así, con el vino y \[A]con el pan
Cris\[G]tos nuevos seamos \[Em]para la hu\[A7]manidad \[Bm-A] 
\endchorus

\beginverse\chordsoff
Sabemos que es largo el camino,
Y lo tenemos que andar
Cristo llévanos de tu mano, Señor,
Para que no nos volvamos atrás.
\endverse

\beginchorus\chordsoff
Y \[D]así, con el vino y \[A]con el pan \dots
\endchorus

\iflink
  \textnote{\href{https://www.youtube.com/watch?v=PYkst40k1_k}{Consultar}} 
\fi
\endsong

% ----------------------------------  Santo ------------------------------------------------- %
\nextcol
\begin{intersong}
\vspace{0.3cm}
\hspace*{-0.7cm}\textbf{\huge\textcolor{darkorange}{--- Santo ---}}
\vspace{0.3cm}
\end{intersong}

% -------------------------------------------------------
% La máxima expresión de alabanza (El grado superlativo)En el idioma hebreo antiguo no existían adjetivos como "santísimo" o "el más santo". Para decir que algo era lo máximo en su categoría, se repetía la palabra tres veces.Decir Kadosh una vez significa: Dios es santo.Decir Kadosh, Kadosh significa: Dios es muy santo.Decir Kadosh, Kadosh, Kadosh significa: Dios es el Santísimo, la perfección absoluta.

\beginsong{Santo Kadosh}[by={}, li={\textbf{Kadosh:} Es una palabra Hebrea que se traduce al español como ``Santo''. Pero su significado es más profundo, es un término que describe la naturaleza pura, la divinidad y el carácter único de Dios.}]

\beginverse
\[Am]Santo, \[C]San\[G]to, \[Am]Santo
\[Am]Santo, \[C]San\[G]to, \[Am]Santo
\endverse

\beginverse
Es el \[Dm]Señor \[F]Dios de los e\[Am]jércitos
Ben\[Dm]dito el que \[F]viene en el \[G]nombre del Se\[Am]ñor
Ho\[Dm]sana en el \[F]cielo, hos\[G]ana al Se\[Am]ñor.
\endverse

\beginverse
\[Am]Kadosh, \[C]Ka\[G]dosh, \[Am]Kadosh.
\[Am]Santo, \[C]San\[G]to, \[Am]Santo
\endverse

\endsong

\begin{intersong}
\vspace{0.3cm}
\hspace*{-0.7cm}\textbf{\huge\textcolor{darkorange}{--- Santo ---}}
\vspace{0.3cm}
\end{intersong}